\documentclass{cskkuproject}

% ==========================================
% ข้อมูลโครงงาน — กรอกข้อมูลจากเอกสาร Proposal
% ==========================================

\titleTH{การวิเคราะห์และทำนายผลการแข่งขันฟุตบอลลีกยุโรปด้วยกระบวนการวิทยาการข้อมูลเพื่อการประเมินมูลค่าการเดิมพัน}
\titleEN{European Football Match Outcome Prediction and Betting Market Analysis Using Data Science}

% ข้อมูลนักศึกษาคนที่ 1
\firstStudentID{683380667-9}
\firstStudentPrefixTH{นาย}
\firstStudentPrefixEN{Mr.}
\firstStudentNameTH{ภาณุพงศ์ ภูสำรอง}
\firstStudentNameEN{Panupong Phoosamrong}

% ข้อมูลนักศึกษาคนที่ 2
\secondStudentID{683380653-0}
\secondStudentPrefixTH{นาย}
\secondStudentPrefixEN{Mr.}
\secondStudentNameTH{ชินภัทร บัวผาย}
\secondStudentNameEN{Chinnapat Buaphai}

% ข้อมูลหลักสูตร
\degreeTH{วิทยาศาสตรบัณฑิต}
\degreeEN{Bachelor of Science}
\majorTH{วิทยาการคอมพิวเตอร์}
\majorEN{Computer Science}

% อาจารย์ที่ปรึกษา
\advisorTH{อ.ดร.ธนพล ตั้งชูพงศ์}
\advisorEN{Lect. Thanapon Tangchoopong, Ph.D.}

% ปีการศึกษา
\academicYearTH{2569}
\academicYearEN{2026}

% รหัสวิชาและชื่อวิชา
\courseCode{CP352101}
\courseName{วิทยาการข้อมูลเบื้องต้น}
\semester{1}

% คำสำคัญ
\keywordsTH{วิทยาการข้อมูล, การทำนายผลการแข่งขันฟุตบอล, ลีกฟุตบอลยุโรป, การเรียนรู้ของเครื่อง, อัตราต่อรอง, การทดสอบย้อนหลัง}
\keywordsEN{Data Science, Football Match Outcome Prediction, European Football Leagues, Machine Learning, Betting Odds, Backtesting}


% ตัวเลขในสมการและตารางใช้ฟอนต์เดียวกับเนื้อความ (Sarabun) ส่วนสัญลักษณ์ใช้ Latin Modern Math
\usepackage{unicode-math}
\setmathfont{Latin Modern Math}
\setmathfont{Sarabun-Regular.ttf}[Path=./theme/fonts/, Scale=1.4, range={up/num, "2E, "2C}]
% หน้าที่มีแต่รูปหรือตาราง วางชิดบนแทนการจัดกลางหน้า
\makeatletter
\setlength{\@fptop}{0pt}
\setlength{\@fpbot}{0pt plus 1fil}
\makeatother
% ชื่อรายการโค้ดเป็นภาษาไทย
\renewcommand{\lstlistingname}{รายการที่}
% หัวข้อยาวตัดบรรทัดแบบชิดซ้าย ไม่ยืดช่องไฟระหว่างคำและไม่ตัดคำด้วยยัติภังค์
\titleformat{\section}{\bfseries\large\upshape\raggedright\hyphenpenalty=10000}{\thesection}{0.45\firstindentlength}{}
\titleformat{\subsection}{\bfseries\large\upshape\raggedright\hyphenpenalty=10000}{\thesubsection}{0.45\firstindentlength}{}
\titleformat{\subsubsection}{\bfseries\large\upshape\raggedright\hyphenpenalty=10000}{\thesubsubsection}{0.45\firstindentlength}{}

\begin{document}

% ==========================================
% หน้าปก
% ==========================================
\maketitlepage

% ตั้งค่าเลขหน้าส่วนต้นเป็นอักษรไทย (ก, ข, ค, ...)
\pagenumbering{arabic}
\setcounter{page}{0}
\renewcommand{\thepage}{\ThaiAlph{\value{page}}}

% ==========================================
% บทคัดย่อ
% ==========================================
\begin{abstractTH}
โครงงานนี้มีวัตถุประสงค์เพื่อประยุกต์ใช้กระบวนการวิทยาการข้อมูลแบบครบวงจรในการวิเคราะห์ปัจจัยก่อนการแข่งขันฟุตบอลลีกยุโรป พัฒนาแบบจำลองทำนายผลการแข่งขัน 3 ทางเลือก (เจ้าบ้านชนะ เสมอ และทีมเยือนชนะ) และตอบคำถามว่าผู้เดิมพันที่ใช้ข้อมูลสาธารณะสามารถเอาชนะตลาดเดิมพันได้หรือไม่

การดำเนินงานได้รวบรวมข้อมูล 5 ลีกใหญ่ของยุโรป ฤดูกาล 2014--2015 ถึงฤดูกาล 2026--2027 ที่กำลังแข่งขันอยู่ จาก Understat, Transfermarkt และ Football-Data.co.uk ผ่านท่อส่งข้อมูลเชิงโมดูลาร์ที่ตรวจสอบโครงสร้างข้อมูลด้วย Pandera แล้วสร้างตัวแปรก่อนแข่ง ได้แก่ คะแนน Elo ดัชนีความเหนื่อยล้า ความหนาแน่นของตารางแข่ง และฟอร์มผู้เล่น 11 ตัวจริง คัดเลือกตัวแปรแบบเดินหน้า และเปรียบเทียบแบบจำลอง Logistic Regression, Gradient-boosted Trees และแบบจำลองจำนวนประตู Poisson โดยทำนายทุกฤดูกาลจากฤดูกาลก่อนหน้าเท่านั้น จากนั้นเปรียบเทียบกับราคาของ Pinnacle ทดสอบกลยุทธ์การเดิมพัน 18 กลยุทธ์ที่ลงทะเบียนไว้ล่วงหน้าพร้อมการปรับแก้แบบ Benjamini--Hochberg และจำลองแผนการวางเงิน 5 แบบ

ผลการศึกษาพบว่า จากข้อมูลพรีเมียร์ลีกฤดูกาล 2000--2001 ถึง 2024--2025 ความได้เปรียบของเจ้าบ้าน ($+0.362$ ประตูต่อนัด) และความสัมพันธ์ของคะแนน Elo กับผลต่างประตู ($r = 0.441$) มีนัยสำคัญทางสถิติ ตัวแปรที่ช่วยทำนายได้คือคะแนน Elo และฟอร์มผู้เล่นตัวจริง ส่วนดัชนีความเหนื่อยล้าไม่ช่วยเมื่อมีคะแนน Elo อยู่แล้ว Logistic Regression ให้ผลดีที่สุด (Log Loss 0.9855, Accuracy 52.41\%) แต่ยังด้อยกว่าราคาเจ้ามือ (0.9691) การวางเดิมพันตามแบบจำลองขาดทุนประมาณร้อยละ 4--5.5 และไม่มีกลยุทธ์ใดใน 252 กรณีที่มีกำไรอย่างมีนัยสำคัญ ความได้เปรียบตามคาดหมายมีเพียงการเลือกราคาดีที่สุดระหว่างเจ้ามือ (Line Shopping) ซึ่งต้องใช้เวลาหลายฤดูกาลกว่าจะเห็นผล และเจ้ามือที่ให้ราคานั้นปิดบัญชีผู้ชนะ ส่วน Martingale และ Kelly เต็มทำให้เงินทุนหมดได้ง่าย ผลเหล่านี้ชี้ว่าตลาดเดิมพันของ 5 ลีกใหญ่มีประสิทธิภาพสูง และท่อส่งข้อมูลที่สร้างซ้ำได้นี้ใช้เป็นแนวทางประเมินกลยุทธ์การเดิมพันกีฬาอย่างเป็นระบบได้

\end{abstractTH}

\begin{abstractEN}
% ABSTRACT (English)
% Note: Should not exceed 300 words.
% Should cover: Objective, Methodology, Results, Applications
% Use Past tense for objectives and methodology
% Use Present tense for results and applications

This project applied an end-to-end data science process to pre-match factors in European football, built models that predict three-way match outcomes (home win, draw, away win), and asked whether a bettor using public data can beat the betting market.

Match data for the five major European leagues from 2014--2015 to the 2026--2027 season in progress was collected from Understat, Transfermarkt and Football-Data.co.uk through a modular pipeline validated with Pandera schemas. Pre-match features were built, including Elo ratings, fatigue indices, schedule congestion and the form of the eleven starters, and selected by forward selection. Logistic regression, gradient-boosted trees and a Poisson goals model were compared, with every season predicted only from the seasons before it. The forecasts were compared with Pinnacle's prices, 18 betting strategies registered before any was run were tested with a Benjamini--Hochberg correction, and five staking plans were simulated.

On Premier League matches from 2000--2001 to 2024--2025, home advantage (+0.362 goals per game) and the link between Elo difference and goal difference ($r = 0.441$) are statistically significant. Elo and the starters' form carry the predictive signal; fatigue adds nothing once Elo is known. Logistic regression performs best (log loss 0.9855, accuracy 52.41\%) but remains behind the bookmaker's prices (0.9691). Betting on the model's positive expected value picks loses 4 to 5.5 percent, and none of the 252 strategy-by-league cells is significant after correction. The only expected edge comes from taking the best price across bookmakers (line shopping), which takes many seasons to show as profit and is offered by bookmakers that close winning accounts. Martingale and full Kelly staking ruin the bankroll easily. The results indicate that the betting markets of the five major leagues are highly efficient, and the reproducible pipeline provides a framework for evaluating sports betting strategies systematically.

\end{abstractEN}

% ==========================================
% กิตติกรรมประกาศ
% ==========================================
\begin{acknowledgement}
% กิตติกรรมประกาศ
% หมายเหตุ: ไม่เกิน 1 หน้ากระดาษ

\emergencystretch=3em
โครงงานฉบับนี้สำเร็จลุล่วงไปได้ด้วยดีด้วยความกรุณาอย่างยิ่งจาก อ.ดร.ธนพล ตั้งชูพงศ์ อาจารย์ที่ปรึกษาโครงงาน ที่ได้ให้คำแนะนำ แนวคิด และข้อเสนอแนะอันมีค่ายิ่ง ตลอดจนคอยตรวจแก้ไขข้อบกพร่องและให้กำลังใจแก่ผู้จัดทำมาโดยตลอด ผู้จัดทำขอกราบขอบพระคุณเป็นอย่างสูงไว้ ณ โอกาสนี้

ขอกราบขอบพระคุณคณาจารย์ประจำสาขาวิชาวิทยาการคอมพิวเตอร์ วิทยาลัยการคอมพิวเตอร์ มหาวิทยาลัยขอนแก่น ทุกท่าน ที่ได้ประสิทธิ์ประสาทวิชาความรู้และถ่ายทอดประสบการณ์ ซึ่งเป็นพื้นฐานสำคัญในการศึกษาค้นคว้าและจัดทำโครงงานฉบับนี้

ขอบคุณเพื่อนร่วมชั้นเรียนและผู้ที่มีส่วนเกี่ยวข้องทุกท่าน ที่คอยให้ความช่วยเหลือ แนะนำ และแลกเปลี่ยนความคิดเห็นตลอดช่วงเวลาของการดำเนินงาน

ท้ายที่สุดนี้ ผู้จัดทำขอกราบขอบพระคุณบิดา มารดา และสมาชิกในครอบครัวทุกท่าน สำหรับความรัก การสนับสนุน และกำลังใจอันยิ่งใหญ่ที่เป็นแรงผลักดันสำคัญให้การศึกษานี้ประสบความสำเร็จ


\end{acknowledgement}

% ==========================================
% สารบัญ
% ==========================================
\tableofcontents
\clearpage
\listoffigures
\clearpage
\listoftables
\clearpage

% ==========================================
% เนื้อเรื่อง 5 บท
% ==========================================
\pagenumbering{arabic}
\setcounter{page}{1}

\chapterone
\chaptertwo
\chapterthree
\chapterfour
\chapterfive

% ==========================================
% เอกสารอ้างอิง
% ==========================================
\clearpage
\addcontentsline{toc}{chapter}{เอกสารอ้างอิง}
\makeatletter
\let\ps@plain\ps@empty
\makeatother
\bibliographystyle{ieeetr}
\bibliography{references}

% ==========================================
% ภาคผนวก
% ==========================================
%\input{appendices}

% ==========================================
% ประวัติผู้เขียน
% ==========================================
%\begin{biography}
%\input{biography}
%\end{biography}

\end{document}